\section{Data Collection and Preprocessing}
\label{sec:data}

\subsection{Dataset}

We collected a custom dataset of hand gesture video clips recorded via
webcam and phone cameras. Each clip captures a single gesture instance
and is labeled as one of five classes: \texttt{grab} (palm $\to$ fist),
\texttt{release} (fist $\to$ palm), \texttt{swipe\_up},
\texttt{swipe\_down}, and \texttt{noise} (invalid/unrecognized motions).
The dataset is split into training and test sets, organized by class.

\subsection{Hand Landmark Extraction}

For each frame of a video clip, we use MediaPipe Hand
Landmarker~\cite{zhang2020mediapipe} to extract 21~3D hand landmarks,
producing a 63-dimensional raw feature vector per frame
($21 \times 3$ coordinates: $x$, $y$, $z$). Coordinates $x$ and $y$
are normalized to $[0, 1]$ relative to the image dimensions, and $z$
represents depth relative to the wrist. If MediaPipe fails to detect a
hand in a given frame, the result is marked as missing. Videos with
fewer than 10\% valid frames are discarded as low-quality.

\subsection{Missing Frame Interpolation}

Because MediaPipe may fail to detect hands in certain frames (due to
motion blur, occlusion, or poor lighting), we apply linear interpolation
to fill gaps between valid detections. For missing frames at the
boundaries of a video, linear extrapolation from the two nearest valid
frames is used. This ensures a complete landmark sequence without
discontinuities.

\subsection{Temporal Resampling}

Since gesture videos vary in duration, we resample every landmark
sequence to a fixed length of $T = 30$~frames using linear
interpolation along the time axis. This standardization ensures
uniform input dimensions for the neural network regardless of the
original video frame rate or gesture speed.

\subsection{Feature Engineering}
\label{subsec:features}

Raw landmark coordinates are not ideal inputs for classification, as they
are sensitive to hand position in the frame and distance from the camera.
We compute a 144-dimensional feature vector per frame from the 63~raw
values, as summarized in Table~\ref{tab:features}.

\begin{table}[ht]
\centering
\caption{Feature engineering: 63 raw landmark values $\to$ 144
  engineered features per frame.}
\label{tab:features}
\begin{tabular}{@{}llp{6.5cm}@{}}
\toprule
\textbf{Feature Group} & \textbf{Dims} & \textbf{Description} \\
\midrule
Normalized landmarks & 63 & Landmarks relative to the wrist, divided by palm size (distance from wrist to middle finger MCP). Translation- and scale-invariant. \\
Velocity & 63 & Frame-to-frame difference of normalized landmarks. Captures motion direction and speed. \\
Wrist velocity & 3 & Absolute wrist position change between frames. Critical for swipe detection. \\
Finger pair distances & 10 & Euclidean distances between all $\binom{5}{2}=10$ fingertip pairs. Measures hand openness. \\
Finger curl angles & 5 & Bending angle (via $\arccos$) for each of the 5~fingers. Distinguishes extended from curled. \\
\midrule
\textbf{Total} & \textbf{144} & \\
\bottomrule
\end{tabular}
\end{table}

\subsection{Data Augmentation}

To expand our limited video dataset, we apply seven augmentation
techniques during preprocessing:

\begin{enumerate}[leftmargin=*]
  \item \textbf{Time crops:} Extract sub-sequences at different
    start/end positions (up to 9 crops per video).
  \item \textbf{Jitter:} Add small Gaussian noise
    ($\sigma = 0.003$) to landmark coordinates (3~variants).
  \item \textbf{Rotation:} Rotate landmarks around the wrist by
    $\{-15, -10, -5, +5, +10, +15\}$~degrees (6~variants).
  \item \textbf{Scale:} Scale hand size by factors
    $\{0.85, 0.9, 1.1, 1.15\}$ relative to the wrist (4~variants).
  \item \textbf{Time warp:} Non-linearly stretch/compress along the
    time axis (2~variants).
  \item \textbf{Speed change:} Uniformly speed up or slow down the
    gesture by a random factor in $[0.8, 1.2]$ (2~variants).
  \item \textbf{Time reversal with label remapping:} Reversing a
    \texttt{grab} produces a valid \texttt{release} and vice versa;
    reversing \texttt{swipe\_up} yields \texttt{swipe\_down}
    (1~variant with corrected label).
\end{enumerate}

This pipeline yields approximately 20--30 augmented samples per
original video. Additionally, online augmentation (50\% jitter,
30\% time warp) is applied stochastically during training for
further regularization.

\subsection{Normalization Statistics}

We compute per-feature mean $\mu_d$ and standard deviation $\sigma_d$
across all training samples and all time steps. During both training
and inference, features are $z$-score normalized:
\begin{equation}
  \hat{x}_d = \frac{x_d - \mu_d}{\sigma_d + \epsilon}, \quad
  \epsilon = 10^{-8}.
  \label{eq:normalize}
\end{equation}
The statistics are serialized alongside the model for deployment
(see Section~\ref{sec:deployment}).
